% Add "answers" option to print the answers
\documentclass[addpoints, answers]{exam}

\usepackage[utf8]{inputenc}
\usepackage{calc}
\usepackage[dvipsnames]{xcolor}
\usepackage{ifthen}
\usepackage{tcolorbox}

% \makebox[\textwidth]{Name:\enspace\hrulefill}
%
% \vspace{5mm}
%\makebox[\textwidth]{Instructor’s name:\enspace\hrulefill}

\definecolor{MyLightGray}{gray}{0.60}

\newcommand{\truefalse}
{		\begin{checkboxes}
			 \choice True
			 \choice False. Explanation:
		\end{checkboxes}
		\vspace{10pt}
}


% Check if 'answers' is an option
\newcommand{\ifanswers}[2]{%
  \ifthenelse{\boolean{exam@answers}}{#1}{#2}%
}



\newcommand{\answerspace}[1]
{
	\ifprintanswers
	\else
		\setlength{\mylength}{10pt}%
		% \vspace{#1\mylength}
		% Make the grey box
		\par\nobreak\noindent
		\makebox[\linewidth][l]{%
			% set inner padding to zero, so only the frame shows
			{\setlength{\fboxsep}{0pt}%
			% draw a color frame, white (or transparent) inside
			\fcolorbox{MyLightGray}{white}{%
			% the inner box has zero width content but fixed height
			\parbox[t][#1\mylength][c]{%
				% subtract twice the frame‐thickness from the full width
				\dimexpr\linewidth-2\fboxrule\relax
			}{\strut}%
			}%
			}%
		}%
		\par
		\vspace{3\mylength}

  	\fi
}
\newlength{\mylength}

\newcommand{\code}[1]{{\texttt{#1}}}


\usepackage{color}
\usepackage{lastpage}
\usepackage{verbatim}
\usepackage{fancyvrb}
\usepackage{graphicx}
\usepackage{float}
\usepackage{minted}
\usepackage{multicol}
\usepackage{amsmath}


\RecustomVerbatimCommand{\VerbatimInput}{VerbatimInput}{}

\newcommand{\Code}[1]{\mintinline{python}{#1}}

\renewcommand{\solutiontitle}{\noindent\textbf{Solution sketch: }}
\definecolor{SolutionColor}{rgb}{0.8,0.9,1}

\cfoot{Page \thepage\ of \pageref{LastPage}}
\rhead{Initials: \rule{20mm}{0.4pt}}

%%%%%%%%%%%%%%%%%%%%%%%%%%%%%%%%%%%%%%%%
\begin{document}
\fvset{numbers=left,numbersep=3pt}
\thispagestyle{empty}

\begin{center}
	{\LARGE CS 408 Fall 2026 -- Symbolic Execution Quiz \\}
\fbox{\fbox{\parbox{5.5in}{\centering
	\begin{itemize}
        \item Try to come up with a solution on your own first, but
          after that talk to your friends and neighbors.

        \item Full points will be rewarded for \textbf{any} (reasonable)
          attempt.
        \end{itemize}
	}}}
\end{center}


\vspace{5mm}
\makebox[\textwidth]{\textbf{Name (please print)}: \rule{100mm}{0.8pt} }

\begin{questions}

\question Consider the following function:

\begin{minted}{python}
def f(a: int, b: int, c: int) -> int:
    x: int = 0
    y: int = 0
    z: int = 0
    if a != 0:
        x = -2
    if b < 5:
        if a == 0 and c != 0:
            y = 1
    z = 2
    return 0
\end{minted}

  \begin{parts}

  \part[5] The symbolic execution tree of this function has six
    leaves. Write down the full set of path constraints at each of
    these nodes, labelling them as \textsf{P1}, \textsf{P2}, \ldots

    \begin{solutionorbox}[3in]
      \begin{itemize}
      \item[P1] $a = 0 \;\wedge\; b \geq 5$
      \item[P2] $a \neq 0 \;\wedge\; b < 5 \;\wedge\; (a = 0 \wedge c \neq 0)$
      \item[P3] $a \neq 0 \;\wedge\; b < 5 \;\wedge\; \neg(a = 0 \wedge c \neq 0)$
      \item[P4]  $a = 0 \;\wedge\; b < 5 \;\wedge\; (a = 0 \wedge c \neq 0)$
      \item[P5]  $a = 0 \;\wedge\; b < 5 \;\wedge\; \neg(a = 0 \wedge
        c \neq 0)$
      \item[P6] $a \neq 0 \;\wedge\; b \geq 5$
      \end{itemize}
    \end{solutionorbox}

    \newpage
  \part[10] For each of the path constraints you identified above,
    either a) provide concrete input values that take the
    corresponding test path or b) identify the path as infeasible.

    \begin{solutionorbox}[3in]
      \begin{itemize}
      \item[P1] \Code{f(0, 5, 10)}
      \item[P2] \textbf{infeasible path}
      \item[P3] \Code{f(1, 4, 42)}
      \item[P4] \Code{f(0, 4, 1)}
      \item[P5] \Code{f(0, 4, 0)}
      \item[P6] \Code{f(4, 6, 0)}
      \end{itemize}
    \end{solutionorbox}

  \part[3] Order the (feasible) paths you identified above by the
    likelihood that three purely random arguments to \Code{f} will
    generate that path, from most to least likely. (No need to be
    perfectly precise here, a rough ranking will do!)

    \begin{solutionorbox}[3in]
      \begin{itemize}
      \item[P3] ($a \neq 0 \wedge b \geq 5$): about \textbf{1 in 2}
      \item[P6] ($a \neq 0 \wedge b \geq 5$): also about \textbf{1 in 2}
      \item[P1] ($a = 0 \wedge b \geq 5$): about \textbf{1 in $2^{33}$}.
      \item[P4] ($a = 0 \wedge b < 5 \wedge c \neq 0$): also about
        \textbf{1 in $2^{33}$}
      \item[P5] ($a = 0 \wedge b < 5 \wedge c = 0$): about \textbf{1 in
          $2^{65}$}, since \emph{two} variables have to be exactly
        zero.
      \end{itemize}
    \end{solutionorbox}

  \part[1] What is your favorite halloween candy?
    \begin{solutionorbox}[1in]
      The correct answer is reeses.
    \end{solutionorbox}
\end{parts}

\end{questions}

\end{document}
